% !TEX root = ../main.tex
\corrigesChapitre{13}{Décrire un mouvement}
\corrigesSouples

\begin{blocLivre}{cor-13-1}
\begin{corrige}{1}{Dériver pour décrire un mouvement}
\demarche{Dériver chaque coordonnée une fois pour la vitesse, deux fois pour l'accélération.}
\begin{questions}
\item $v_x=\ddt{x}=\res{3{,}0}$ et $v_y=\ddt{y}=\res{-4{,}0\,t+4{,}0}$ ; puis $a_x=\res{0}$ et $a_y=\res{-4{,}0}$ (en unités SI).
\item $\vec{a}$ est \res{\text{constant}} : direction verticale, vers le bas, norme \qty{4,0}{m.s^{-2}}. Le mouvement n'est
pas rectiligne : $\vec{v}$ n'est pas colinéaire à $\vec{a}$ (sa coordonnée $v_x$ est non nulle).
\item À $t=0$ : $v=\sqrt{3{,}0^2+4{,}0^2}=\res{\qty{5,0}{m.s^{-1}}}$ ; à $t=\qty{1,0}{s}$ : $v_y=0$ et
$v=\res{\qty{3,0}{m.s^{-1}}}$. Comme $v_x$ est constant, $v$ est minimale quand $v_y=0$, donc à
$\res{t=\qty{1,0}{s}}$.
\item $t=\frac{x-1{,}0}{3{,}0}$, d'où $\res{y=-\frac{2{,}0}{9}(x-1{,}0)^2+\frac{4{,}0}{3}(x-1{,}0)}$ : la trajectoire est un
arc de parabole.
\end{questions}
\end{corrige}
\end{blocLivre}
\begin{blocLivre}{cor-13-2}
\begin{corrige}{2}{Trois chronophotographies}
\demarche{Comparer les distances parcourues pendant des durées égales ; la vitesse en $M_i$ vaut $M_{i-1}M_{i+1}/(2\tau)$.}
\begin{questions}
\item A : distances égales, \res{\text{rectiligne uniforme}} ; B : distances croissantes, \res{\text{rectiligne
accéléré}} ; C : distances décroissantes, \res{\text{rectiligne ralenti}}.
\item B : $\vec{a}$ \res{\text{de même sens}} que $\vec{v}$ ; C : $\vec{a}$ \res{\text{de sens contraire}}.
\item $v_\mathrm{A}=\dfrac{0{,}60}{0{,}20}=\res{\qty{3,0}{m.s^{-1}}}$.
\item $v_2=\dfrac{0{,}40+0{,}60}{2\times0{,}20}=\res{\qty{2,5}{m.s^{-1}}}$ et
$v_4=\dfrac{0{,}80+1{,}00}{0{,}40}=\res{\qty{4,5}{m.s^{-1}}}$, d'où
$a_3=\dfrac{v_4-v_2}{2\tau}=\dfrac{2{,}0}{0{,}40}=\res{\qty{5,0}{m.s^{-2}}}$.
\end{questions}
\end{corrige}
\end{blocLivre}
\begin{blocLivre}{cor-13-3}
\begin{corrige}{3}{L'accélération du physicien}
\demarche{L'accélération est non nulle dès que le vecteur vitesse change, en norme ou en direction.}
\begin{questions}
\item {\renewcommand{\arraystretch}{1.25}\begin{tabular}[t]{@{}l@{\hspace{3mm}}l@{}}
a. & vecteur vitesse constant : $\res{\vec{a}=\vec{0}}$ ;\\
b. & la direction de $\vec{v}$ change : $\vec{a}$ \res{\text{dirigé vers le centre du virage}} ;\\
c. & $\vec{a}$ colinéaire à la route, \res{\text{opposé à la vitesse}} ;\\
d. & $\vec{a}$ \res{\text{vertical, vers le bas}} ;\\
e. & $\vec{a}$ \res{\text{dirigé vers le centre de la Terre}} ;\\
f. & vecteur vitesse constant : $\res{\vec{a}=\vec{0}}$.
\end{tabular}}
\item La valeur de la vitesse est constante, mais sa direction change : le vecteur vitesse varie, l'accélération
\res{\text{n'est pas nulle}} ; elle vaut $v^2/R$ et est dirigée vers le centre du virage.
\end{questions}
\end{corrige}
\end{blocLivre}
\begin{blocLivre}{cor-13-4}
\begin{corrige}{4}{Les chaises volantes}
\demarche{Mouvement circulaire uniforme : $v=2\pi R/T$ et $a=v^2/R$, vers le centre.}
\begin{questions}
\item La nacelle décrit un cercle et fait chaque tour dans la même durée : la valeur de sa vitesse est constante, le
mouvement est \res{\text{circulaire uniforme}}.
\item $v=\dfrac{2\pi R}{T}=\dfrac{2\pi\times8{,}0}{6{,}0}=\res{\qty{8,4}{m.s^{-1}}}$.
\item $\vec{a}$ est horizontal, dirigé vers l'axe du manège (centripète), de norme
$a=\dfrac{v^2}{R}=\dfrac{8{,}38^2}{8{,}0}=\res{\qty{8,8}{m.s^{-2}}}$, soit environ $\res{0{,}9\,g}$.
\item $\ut$ tangent au cercle dans le sens de rotation, $\un$ vers l'axe ; $\vec{v}$ porté par $\ut$, $\vec{a}$ porté par
$\un$.
\end{questions}
\schema[corrige]{rev-13-ex04-cor}
\legende{Vue de dessus. Les longueurs des flèches de vitesse et d’accélération utilisent des échelles distinctes.}
\end{corrige}
\end{blocLivre}
\begin{blocLivre}{cor-13-5}
\begin{corrige}{5}{Lire un vecteur accélération}
\demarche{La coordonnée tangentielle indique si la vitesse augmente ($a_\mathrm{t}>0$) ou diminue ($a_\mathrm{t}<0$).}
\begin{questions}
\item En A, $\vec{a}$ est dirigé vers le centre : $a_\mathrm{t}=0$, vitesse \res{\text{constante}}. En B, $\vec{a}$ penche
vers l'avant : $a_\mathrm{t}>0$, vitesse \res{\text{croissante}}. En C, il penche vers l'arrière : $a_\mathrm{t}<0$,
vitesse \res{\text{décroissante}}.
\item $a_\mathrm{n}=\dfrac{v^2}{R}=\dfrac{12^2}{50}=\res{\qty{2,9}{m.s^{-2}}}$ et
$a=\sqrt{1{,}5^2+2{,}88^2}=\res{\qty{3,2}{m.s^{-2}}}$.
\item En ligne droite à vitesse constante, $\res{\vec{a}=\vec{0}}$.
\end{questions}
\end{corrige}
\end{blocLivre}
\begin{blocLivre}{cor-13-6}
\begin{corrige}{6}{Le départ d'un sprinteur}
\demarche{Phase 1 : uniformément accélérée ; phase 2 : uniforme. La distance est l'aire sous la courbe $v(t)$.}
\begin{questions}
\item Phase 1 : \res{\text{rectiligne uniformément accéléré}} ; phase 2 : \res{\text{rectiligne uniforme}}.
\item $a=\dfrac{11{,}6}{4{,}0}=\res{\qty{2,9}{m.s^{-2}}}$ ; $v(t)=2{,}9\,t$ et $x(t)=1{,}45\,t^2$.
\item $x(4{,}0)=1{,}45\times4{,}0^2=\res{\qty{23,2}{m}}$ ; l'aire du triangle sous la courbe vaut
$\frac12\times4{,}0\times11{,}6=\qty{23,2}{m}$.
\item Il reste \qty{76,8}{m} à \qty{11,6}{m.s^{-1}}, soit \qty{6,6}{s} : temps total \res{\qty{10,6}{s}}.
\end{questions}
\schema[corrige]{13-ex06}
\end{corrige}
\end{blocLivre}
\begin{blocLivre}{cor-13-7}
\begin{corrige}{7}{La distance d'arrêt}
\demarche{Mouvement uniforme pendant la réaction, puis uniformément ralenti pendant le freinage.}
\begin{questions}
\item $v_0=\qty{25}{m.s^{-1}}$ : $d_\mathrm{R}=25\times1{,}0=\res{\qty{25}{m}}$, mouvement \res{\text{rectiligne uniforme}}.
\item Avec un axe dans le sens du mouvement : $v(t)=25-a\,t$ et $x(t)=25\,t-\frac12\,a\,t^2$.
\item $v=0$ pour $t_\mathrm{F}=\dfrac{25}{7{,}5}=\res{\qty{3,3}{s}}$ ; $d_\mathrm{F}=25\,t_\mathrm{F}-\frac12\times7{,}5\,t_\mathrm{F}^2=\res{\qty{42}{m}}$ ;
distance d'arrêt : $25+42=\res{\qty{67}{m}}$.
\item $d_\mathrm{F}=\dfrac{v_0^2}{2a}=\dfrac{625}{8{,}0}=\qty{78}{m}$ : distance d'arrêt \res{\qty{103}{m}}. Sur route mouillée,
la distance d'arrêt augmente de moitié : il faut réduire la vitesse.
\end{questions}
\end{corrige}
\end{blocLivre}
\begin{blocLivre}{cor-13-8}
\begin{corrige}{8}{Une chute filmée}
\demarche{Les vitesses croissent régulièrement : le coefficient directeur de $v_y(t)$ est l'accélération.}
\begin{questions}
\item $v_i\approx\dfrac{y_{i+1}-y_{i-1}}{2\tau}$ : vitesse moyenne entre les dates encadrantes. Les cinq vitesses
correspondent aux dates \res{\num{0,050}, \num{0,10}, \num{0,15}, \num{0,20}\text{ et }\qty{0,25}{s}}.
\item Les points sont alignés sur une droite passant par l'origine : \res{\text{mouvement rectiligne uniformément
accéléré}}.
\item Les vitesses augmentent de \qty{0,49}{m.s^{-1}} toutes les \qty{0,050}{s} :
$a=\dfrac{0{,}49}{0{,}050}=\res{\qty{9,8}{m.s^{-2}}}$, la valeur de l'intensité de la pesanteur.
\item \code{a = [(v[i+1] - v[i-1]) / (2 * tau) for i in range(1, 4)]} puis \code{print(a)}.
\end{questions}
\schema[corrige]{rev-13-ex08-cor}
\legende{La pente de la droite vaut $g$ pour l’axe vertical orienté vers le bas.}
\end{corrige}
\end{blocLivre}
\begin{blocLivre}{cor-13-9}
\begin{corrige}{9}{Sur un rond-point}
\demarche{$a_\mathrm{t}=\Delta v/\Delta t$ (constante) et $a_\mathrm{n}=v^2/R$ ; les deux coordonnées se composent.}
\begin{questions}
\item $a_\mathrm{t}=\dfrac{10{,}0-5{,}0}{4{,}0}=\res{\qty{1,25}{m.s^{-2}}}$.
\item $a_\mathrm{n}=\dfrac{8{,}0^2}{20}=\res{\qty{3,2}{m.s^{-2}}}$ et $a=\sqrt{1{,}25^2+3{,}2^2}=\res{\qty{3,4}{m.s^{-2}}}$.
\item $\vec{v}$ porté par $\ut$ ; $\vec{a}$ dirigé vers l'intérieur et vers l'avant : $\tan\theta=\dfrac{1{,}25}{3{,}2}$,
soit $\res{\theta=21^\circ}$ avec $\un$.
\item $a_\mathrm{n}$ \res{\text{augmente}} comme le carré de la vitesse : de \qty{1,25}{m.s^{-2}} à
\qty{5,0}{m.s^{-2}}.
\end{questions}
\schema[corrige]{rev-13-ex09-cor}
\legende{L’accélération est la somme de ses composantes tangentielle et normale.}
\end{corrige}
\end{blocLivre}
\begin{blocLivre}{cor-13-10}
\begin{corrige}{10}{La Station spatiale internationale}
\demarche{Mouvement circulaire uniforme dans le référentiel géocentrique : $v=2\pi r/T$ et $a=v^2/r$.}
\begin{questions}
\item $r=6371+420=\qty{6791}{km}$ et $v=\dfrac{2\pi\times\num{6.791e6}}{92{,}8\times60}=\res{\qty{7,66e3}{m.s^{-1}}}$, soit
\res{\num{27600}~\text{km/h}}.
\item $\vec{a}$ dirigé vers le centre de la Terre, de norme $a=\dfrac{v^2}{r}=\res{\qty{8,65}{m.s^{-2}}}$.
\item $\dfrac{24\times60}{92{,}8}=\res{15{,}5}$ tours par jour.
\item $a$ vaut 88~\% de $g$ au sol : à 420~km d'altitude, l'attraction de la Terre reste forte. La station « tombe »
en permanence autour de la Terre.
\end{questions}
\end{corrige}
\end{blocLivre}
\begin{blocLivre}{cor-13-11}
\begin{corrige}{11}{Représenter l'accélération d'un mouvement circulaire}
\demarche{La différence seconde divisée par $\tau^2$ approche la dérivée seconde de la position.}
\begin{questions}
\item $\theta=\dfrac{2\pi\times0{,}5}{4{,}0}=\dfrac{\pi}{4}$ : $x=y=2{,}0\cos\dfrac{\pi}{4}=\res{\qty{1,41}{m}}$.
\item $\dfrac{x(t+\tau)-x(t)}{\tau}$ approche $v_x$ à $t+\frac{\tau}{2}$ et $\dfrac{x(t)-x(t-\tau)}{\tau}$ approche $v_x$ à
$t-\frac{\tau}{2}$ ; leur différence divisée par $\tau$ approche $\ddt{v_x}=a_x$ à la date $t$.
\item $\vec{a}$ a pour coordonnées $(-3{,}48\,;\,-3{,}48)$ : il est porté par la droite MO, \res{\text{dirigé vers le
centre du cercle}}.
\item $a=\sqrt{2}\times3{,}482=\res{\qty{4,92}{m.s^{-2}}}$, contre $\dfrac{v^2}{R}=R\left(\dfrac{2\pi}{T}\right)^2=\qty{4,93}{m.s^{-2}}$ :
écart inférieur à 0,2~\%.
\end{questions}
\schema[corrige]{rev-13-ex11-cor}
\legende{La différence seconde donne un vecteur dirigé vers le centre du cercle.}
\end{corrige}
\end{blocLivre}
\begin{blocLivre}{cor-13-12}
\begin{corrige}{12}{La centrifugeuse des pilotes}
\demarche{En régime établi, $a=v^2/R$ ; pendant la montée en régime, $a_\mathrm{t}$ et $a_\mathrm{n}$ coexistent.}
\begin{questions}
\item \res{\text{Circulaire uniforme}} : $\vec{a}$ horizontal, dirigé vers l'axe de rotation, de norme constante
$v^2/R$.
\item $v=\sqrt{8g\,R}=\sqrt{78{,}5\times7{,}6}=\res{\qty{24}{m.s^{-1}}}$ ; fréquence
$f=\dfrac{v}{2\pi R}=\qty{0,51}{tr.s^{-1}}$, soit \res{31~\text{tours par minute}}.
\item $a_\mathrm{t}=\dfrac{24{,}4}{20}=\res{\qty{1,2}{m.s^{-2}}}$.
\item À mi-parcours, $v=\qty{12,2}{m.s^{-1}}$ et $a_\mathrm{n}=\dfrac{12{,}2^2}{7{,}6}=\qty{19,6}{m.s^{-2}}$ ($2\,g$) :
$a=\sqrt{19{,}6^2+1{,}2^2}=\res{\qty{19,7}{m.s^{-2}}}$, incliné de \res{3,6^\circ} seulement par rapport au bras : la
composante normale domine.
\end{questions}
\end{corrige}
\end{blocLivre}
\begin{blocLivre}{cor-13-13}
\begin{corrige}{13}{Le lancer franc}
\demarche{Dériver les équations horaires, puis exploiter $v_y=0$ au sommet et $x=\qty{4,60}{m}$ au panier.}
\begin{questions}
\item En dérivant : $v_x=4{,}78$ et $v_y=-9{,}80\,t+5{,}70$, puis $a_x=0$ et $a_y=-9{,}80$. Le vecteur accélération est
\res{\text{constant, vertical, vers le bas}}, de norme \qty{9,80}{m.s^{-2}}.
\item $v_0=\sqrt{4{,}78^2+5{,}70^2}=\res{\qty{7,44}{m.s^{-1}}}$ et $\tan\alpha=\dfrac{5{,}70}{4{,}78}$, soit $\res{\alpha=50^\circ}$.
\item $v_y=0$ pour $t_\mathrm{S}=\dfrac{5{,}70}{9{,}80}=\res{\qty{0,582}{s}}$ ; $y(t_\mathrm{S})=\res{\qty{3,76}{m}}$.
\item $x=\qty{4,60}{m}$ pour $t=\dfrac{4{,}60}{4{,}78}=\qty{0,962}{s}$ ; alors $y=\qty{3,05}{m}$ : \res{\text{le ballon passe par
le centre de l'anneau}}, en descendant.
\end{questions}
\end{corrige}
\end{blocLivre}
\begin{blocLivre}{cor-13-14}
\begin{corrige}{14}{Le looping d'un grand huit}
\demarche{Au voisinage d'un point, $a_\mathrm{n}=v^2/r$ où $r$ est le rayon de courbure local.}
\begin{questions}
\item En A : $a_\mathrm{n}=\dfrac{22^2}{15}=\res{\qty{32}{m.s^{-2}}}$, soit $\res{3{,}3\,g}$ ; en B :
$a_\mathrm{n}=\dfrac{12^2}{8{,}0}=\res{\qty{18}{m.s^{-2}}}$, soit $\res{1{,}8\,g}$.
\item En A, $\vec{a}$ est vertical, \res{\text{vers le haut}} (vers le centre de courbure) ; en B, vertical, \res{\text{vers
le bas}}.
\item $a=\dfrac{22^2}{8{,}0}=\qty{60}{m.s^{-2}}$, soit \res{6{,}2\,g} : au-delà de la limite de $5\,g$. Le grand rayon en bas
limite l'accélération là où la vitesse est maximale ; le petit rayon en haut permet de passer le sommet à vitesse
réduite.
\end{questions}
\schema[corrige]{13-ex14}
\end{corrige}
\end{blocLivre}
\begin{blocLivre}{cor-13-15}
\begin{corrige}{15}{Le TGV dans les courbes}
\demarche{Convertir les vitesses en \unit{m.s^{-1}} ; $a_\mathrm{n}=v^2/R$ ; freinage : $v_0^2=2ad$.}
\begin{questions}
\item $v=\dfrac{320}{3{,}6}=\qty{88,9}{m.s^{-1}}$ et $a_\mathrm{n}=\dfrac{88{,}9^2}{7000}=\res{\qty{1,1}{m.s^{-2}}}$.
\item $v=\sqrt{1{,}13\times1200}=\qty{37}{m.s^{-1}}$, soit \res{\text{environ 130~km/h}} : sur ligne classique, le TGV doit
ralentir fortement dans les courbes serrées.
\item $a=\dfrac{v_0^2}{2d}=\dfrac{88{,}9^2}{2\times3300}=\res{\qty{1,2}{m.s^{-2}}}$ et $t=\dfrac{v_0}{a}=\res{\qty{74}{s}}$.
\item Les deux accélérations valent environ \res{0{,}12\,g} : elles restent faibles parce que les rayons des courbes et
les distances de freinage sont très grands.
\end{questions}
\end{corrige}
\end{blocLivre}
